\documentclass[12pt]{article}

\usepackage[utf8]{inputenc}
\usepackage{amsmath, amssymb, amsfonts, cancel}
\usepackage{graphicx}
\usepackage[letterpaper, left=1in, right=1in, top=1in, bottom=1in]{geometry}
\usepackage{booktabs}
\usepackage{multirow}
\usepackage{float}
\usepackage{caption}
\usepackage{colortbl}
\usepackage{xcolor}
\definecolor{paleYellow}{RGB}{255,255,180}
\usepackage{physics}
\usepackage{hyperref}
\usepackage[spanish]{babel}

\title{
Práctica 1: Evaluación Integral de Diseño de una Red Segmentada\\[0.5em]
\large Interacción de redes
}

\author{
Jorge Alberto Suarez Saldana\\
\small Matrícula: 355992\\
\small Profesor: Raymundo Antonio González Grimaldo
}

\date{\today}

\begin{document}
\maketitle

\section{Topología de la red}

Podemos observar en la Figura \ref{fig:topologia} un diagrama de la topología
de la red implementada en el simulador Cisco Packet Tracer. La topología está
compuesta por tres routers, tres switches, diferentes hosts y
un servidor.

Los routers R1, R2 y R3 se encuentran conectados mediante enlaces seriales,
mientras que los switches permiten conectar los dispositivos finales
correspondientes a las diferentes VLAN de la red. Para permitir la comunicación
entre las distintas VLAN se utilizó el método de \textit{router-on-a-stick},
mediante subinterfaces en los routers R2 y R3.

\begin{figure}[H]
  \centering
  \includegraphics[width=0.9\textwidth]{topologia.png}
  \caption{Topología de la red}
  \label{fig:topologia}
\end{figure}


\section{Tabla de direccionamiento}

La red fue dividida en diferentes subredes de acuerdo con la función de cada
segmento de la topología. Para las conexiones entre routers se utilizaron
redes con prefijo /30, mientras que para las redes LAN y las VLAN se utilizaron
subredes /24.

\bigskip

\begin{center}
\renewcommand{\arraystretch}{1.2}
\begin{tabular}{|l|l|l|l|l|}
\hline
\textbf{Dispositivo} & \textbf{VLAN} & \textbf{IP} &
\textbf{Máscara de red} & \textbf{Gateway} \\
\hline

\multicolumn{5}{|c|}{\textbf{Red entre R1 y R2}} \\
\hline
R1 & -- & 10.10.10.1 & 255.255.255.252 & -- \\
R2 & -- & 10.10.10.2 & 255.255.255.252 & -- \\
\hline

\multicolumn{5}{|c|}{\textbf{Red PC-A}} \\
\hline
R1 & -- & 148.120.1.1 & 255.255.255.0 & -- \\
PC-A & -- & 148.120.1.10 & 255.255.255.0 & 148.120.1.1 \\
\hline

\multicolumn{5}{|c|}{\textbf{Red PC-B}} \\
\hline
R2 & -- & 148.120.2.1 & 255.255.255.0 & -- \\
PC-B & -- & 148.120.2.10 & 255.255.255.0 & 148.120.2.1 \\
\hline

\multicolumn{5}{|c|}{\textbf{Enlace entre R2 y R3}} \\
\hline
R2 & -- & 10.10.20.1 & 255.255.255.252 & -- \\
R3 & -- & 10.10.20.2 & 255.255.255.252 & -- \\
\hline

\multicolumn{5}{|c|}{\textbf{VLAN 20 - Nativa}} \\
\hline
R2 & 20 & 192.168.20.1 & 255.255.255.0 & -- \\
S1 & 20 & 192.168.20.2 & 255.255.255.0 & 192.168.20.1 \\
S2 & 20 & 192.168.20.3 & 255.255.255.0 & 192.168.20.1 \\
\hline

\multicolumn{5}{|c|}{\textbf{VLAN 25 - Coordinadores}} \\
\hline
R2 & 25 & 192.168.25.1 & 255.255.255.0 & -- \\
PC-D & 25 & 192.168.25.10 & 255.255.255.0 & 192.168.25.1 \\
LAPTOP-B & 25 & 192.168.25.11 & 255.255.255.0 & 192.168.25.1 \\
\hline

\multicolumn{5}{|c|}{\textbf{VLAN 30 - Profesores}} \\
\hline
R2 & 30 & 192.168.30.1 & 255.255.255.0 & -- \\
LAPTOP-A & 30 & 192.168.30.10 & 255.255.255.0 & 192.168.30.1 \\
PC-C & 30 & 192.168.30.11 & 255.255.255.0 & 192.168.30.1 \\
\hline

\multicolumn{5}{|c|}{\textbf{VLAN 10 - Becarios}} \\
\hline
R3 & 10 & 172.16.10.1 & 255.255.255.0 & -- \\
PC-E & 10 & 172.16.10.10 & 255.255.255.0 & 172.16.10.1 \\
\hline

\multicolumn{5}{|c|}{\textbf{VLAN 20 - Nativa en S3}} \\
\hline
R3 & 20 & 172.16.20.1 & 255.255.255.0 & -- \\
S3 & 20 & 172.16.20.2 & 255.255.255.0 & 172.16.20.1 \\
\hline

\multicolumn{5}{|c|}{\textbf{Red del servidor}} \\
\hline
R3 & -- & 202.20.3.1 & 255.255.255.0 & -- \\
Server & -- & 202.20.3.10 & 255.255.255.0 & 202.20.3.1 \\
\hline

\end{tabular}
\end{center}


\section{Configuración de VLANs}

En los switches S1 y S2 se configuraron tres VLANs: la VLAN 20, utilizada
como VLAN nativa y de administración; la VLAN 25, correspondiente a los
Coordinadores; y la VLAN 30, correspondiente a los Profesores.

Por otro lado, en el switch S3 se configuraron la VLAN 20 como VLAN nativa y
la VLAN 10 para los Becarios.

Los puertos que conectan dispositivos finales fueron configurados como
puertos de acceso pertenecientes a la VLAN correspondiente, mientras que los
enlaces entre switches y routers fueron configurados como enlaces troncales
(\textit{trunk}).

Además, se implementó seguridad de puertos en los puertos de acceso mediante
\textit{port-security}, limitando a un dispositivo la cantidad máxima de
direcciones MAC permitidas.


\section{Configuración de los routers}

Para permitir la comunicación entre las diferentes VLANs se utilizó la
configuración conocida como \textit{router-on-a-stick}. Este método consiste
en utilizar una sola interfaz física del router y dividirla lógicamente
mediante diferentes subinterfaces, asignando a cada una una VLAN mediante
encapsulación IEEE 802.1Q.

En R2 se utilizaron las subinterfaces correspondientes a las VLAN 20, 25 y
30. La VLAN 20 fue configurada como VLAN nativa.

De manera similar, en R3 se utilizaron las subinterfaces correspondientes a
las VLAN 20 y 10, utilizando nuevamente la VLAN 20 como VLAN nativa.

Los enlaces seriales entre los routers fueron configurados utilizando las
redes 10.10.10.0/30 y 10.10.20.0/30. En el extremo DCE de cada enlace se
configuró una velocidad de reloj de 64000 bits por segundo mediante el
comando \textit{clock rate}.


\section{Enrutamiento estático}

Para permitir la comunicación entre todas las redes de la topología se
configuraron rutas estáticas en los tres routers.

R1 utiliza a R2 como siguiente salto para alcanzar las redes que se encuentran
detrás de este router. De manera similar, R2 utiliza a R3 como siguiente salto
para las redes ubicadas en el extremo derecho de la topología.

Finalmente, R3 utiliza a R2 como siguiente salto para alcanzar las redes
ubicadas detrás de R1 y las VLAN configuradas en R2.

Esta configuración permite que un paquete enviado desde una red pueda
atravesar los diferentes routers hasta llegar a la red destino.


\section{Configuración actual de S1}

Para comprobar la configuración activa del switch S1 se utilizó el comando
\textit{show running-config}. A continuación se deja el espacio para colocar
la salida obtenida directamente desde Cisco Packet Tracer.

\begin{verbatim}

S1#show running-config
Building configuration...

Current configuration : 1592 bytes
!
version 15.0
no service timestamps log datetime msec
no service timestamps debug datetime msec
service password-encryption
!
hostname S1
!
enable secret 5 $1$mERr$9cTjUIEqNGurQiFU.ZeCi1
!
!
!
no ip domain-lookup
!
!
!
spanning-tree mode pvst
spanning-tree extend system-id
!
interface FastEthernet0/1
!
interface FastEthernet0/2
!
interface FastEthernet0/3
!
interface FastEthernet0/4
!
interface FastEthernet0/5
 switchport trunk native vlan 20
 switchport mode trunk
!
interface FastEthernet0/6
 switchport access vlan 30
 switchport mode access
 switchport port-security
!
interface FastEthernet0/7
 switchport access vlan 25
 switchport mode access
 switchport port-security
!
interface FastEthernet0/8
!
interface FastEthernet0/9
!
interface FastEthernet0/10
!
interface FastEthernet0/11
!
interface FastEthernet0/12
!
interface FastEthernet0/13
!
interface FastEthernet0/14
!
interface FastEthernet0/15
!
interface FastEthernet0/16
!
interface FastEthernet0/17
!
interface FastEthernet0/18
!
interface FastEthernet0/19
!
interface FastEthernet0/20
!
interface FastEthernet0/21
!
interface FastEthernet0/22
!
interface FastEthernet0/23
!
interface FastEthernet0/24
!
interface GigabitEthernet0/1
 switchport trunk native vlan 20
 switchport mode trunk
!
interface GigabitEthernet0/2
!
interface Vlan1
 no ip address
 shutdown
!
interface Vlan20
 ip address 192.168.20.2 255.255.255.0
!
ip default-gateway 192.168.20.1
!
banner motd ^CACCESO RESTRINGIDO - SOLO PERSONAL AUTORIZADO^C
!
!
!
line con 0
 password 7 0822455D0A16
 login
!
line vty 0 4
 login
line vty 5 15
 login
!
!
!
!
end


S1#
\end{verbatim}


\section{Configuración actual de R2}

De igual manera, se verificó la configuración activa del router R2 mediante
el comando \textit{show running-config}. Esta configuración permite comprobar
las interfaces, subinterfaces, VLANs, direccionamiento y demás parámetros
establecidos durante la práctica.

\begin{verbatim}
R2#show running-config
Building configuration...

Current configuration : 1620 bytes
!
version 15.1
no service timestamps log datetime msec
no service timestamps debug datetime msec
service password-encryption
!
hostname R2
!
!
!
enable secret 5 $1$mERr$9cTjUIEqNGurQiFU.ZeCi1
!
!
!
!
!
!
no ip cef
no ipv6 cef
!
!
!
!
license udi pid CISCO2911/K9 sn FTX15246G8Z-
!
!
!
!
!
!
!
!
!
!
!
no ip domain-lookup
!
!
spanning-tree mode pvst
!
!
!
!
!
!
interface GigabitEthernet0/0
 ip address 148.120.2.1 255.255.255.0
 duplex auto
 speed auto
!
interface GigabitEthernet0/1
 no ip address
 duplex auto
 speed auto
!
interface GigabitEthernet0/1.20
 encapsulation dot1Q 20 native
 ip address 192.168.20.1 255.255.255.0
!
interface GigabitEthernet0/1.25
 encapsulation dot1Q 25
 ip address 192.168.25.1 255.255.255.0
!
interface GigabitEthernet0/1.30
 encapsulation dot1Q 30
 ip address 192.168.30.1 255.255.255.0
!
interface GigabitEthernet0/2
 no ip address
 duplex auto
 speed auto
 shutdown
!
interface Serial0/0/0
 no ip address
 clock rate 2000000
 shutdown
!
interface Serial0/0/1
 no ip address
 clock rate 2000000
 shutdown
!
interface Serial0/1/0
 ip address 10.10.10.2 255.255.255.252
!
interface Serial0/1/1
 ip address 10.10.20.1 255.255.255.252
 clock rate 64000
!
interface Vlan1
 no ip address
 shutdown
!
ip classless
ip route 148.120.1.0 255.255.255.0 10.10.10.1 
ip route 172.16.10.0 255.255.255.0 10.10.20.2 
ip route 172.16.20.0 255.255.255.0 10.10.20.2 
ip route 202.20.3.0 255.255.255.0 10.10.20.2 
!
ip flow-export version 9
!
!
!
banner motd ^CACCESO RESTRINGIDO - SOLO PERSONAL AUTORIZADO^C
!
!
!
!
line con 0
 password 7 0822455D0A16
 login
!
line aux 0
!
line vty 0 4
 login
!
!
!
end


R2#
\end{verbatim}


\section{Tabla de enrutamiento de R1}

Para comprobar las rutas conocidas por R1 se utilizó el comando
\textit{show ip route}. La tabla permite observar tanto las redes directamente
conectadas como las rutas estáticas configuradas para alcanzar las redes
remotas.

\begin{verbatim}
R1#show ip route
Codes: L - local, C - connected, S - static, R - RIP, M - mobile, B - BGP
       D - EIGRP, EX - EIGRP external, O - OSPF, IA - OSPF inter area
       N1 - OSPF NSSA external type 1, N2 - OSPF NSSA external type 2
       E1 - OSPF external type 1, E2 - OSPF external type 2, E - EGP
       i - IS-IS, L1 - IS-IS level-1, L2 - IS-IS level-2, ia - IS-IS inter area
       * - candidate default, U - per-user static route, o - ODR
       P - periodic downloaded static route

Gateway of last resort is not set

     10.0.0.0/8 is variably subnetted, 3 subnets, 2 masks
C       10.10.10.0/30 is directly connected, Serial0/1/0
L       10.10.10.1/32 is directly connected, Serial0/1/0
S       10.10.20.0/30 [1/0] via 10.10.10.2
     148.120.0.0/16 is variably subnetted, 3 subnets, 2 masks
C       148.120.1.0/24 is directly connected, GigabitEthernet0/1
L       148.120.1.1/32 is directly connected, GigabitEthernet0/1
S       148.120.2.0/24 [1/0] via 10.10.10.2
     172.16.0.0/24 is subnetted, 2 subnets
S       172.16.10.0/24 [1/0] via 10.10.10.2
S       172.16.20.0/24 [1/0] via 10.10.10.2
S    192.168.20.0/24 [1/0] via 10.10.10.2
S    192.168.25.0/24 [1/0] via 10.10.10.2
S    192.168.30.0/24 [1/0] via 10.10.10.2
S    202.20.3.0/24 [1/0] via 10.10.10.2

R1#
\end{verbatim}


\section{Pruebas de conectividad}

Una vez terminada la configuración de la topología se realizaron diferentes
pruebas de conectividad mediante el comando \textit{ping}. Estas pruebas
permitieron comprobar la comunicación tanto dentro de una misma VLAN como
entre diferentes redes.

Entre las pruebas realizadas se encuentran la comunicación entre PC-A y
PC-B, LAPTOP-B y PC-D, PC-C y PC-E, R1 y el servidor, así como la comunicación
entre LAPTOP-A y la interfaz de administración del switch S3.

En la Figura \ref{fig:pings} se muestran las pruebas de conectividad
realizadas durante la práctica.

\begin{figure}[H]
  \centering
  \includegraphics[width=0.9\textwidth]{pings.png}
  \caption{Pruebas de conectividad entre diferentes dispositivos}
  \label{fig:pings}
\end{figure}


\section{Preguntas}

\subsection{¿Qué nombre se le da a la configuración de una sola interfaz física
de un router que se utiliza para enrutar el tráfico entre múltiples VLANs
en una red?}

Esta configuración recibe el nombre de \textit{router-on-a-stick}. Consiste
en utilizar una única interfaz física del router dividida en varias
subinterfaces lógicas. Cada subinterfaz se asocia con una VLAN diferente
mediante el protocolo de encapsulación IEEE 802.1Q.

De esta manera, el router puede recibir y enviar tráfico perteneciente a
múltiples VLANs utilizando un único enlace físico configurado como
\textit{trunk} entre el switch y el router.


\subsection{¿Por qué es importante configurar la seguridad de los puertos
de un switch?}

La seguridad de puertos es importante porque permite controlar qué
dispositivos pueden conectarse físicamente a determinados puertos del
switch. Mediante \textit{port-security} es posible limitar la cantidad de
direcciones MAC permitidas y establecer una acción en caso de que se detecte
un dispositivo no autorizado.

Esto ayuda a prevenir que una persona conecte un dispositivo no autorizado
a un puerto disponible de la red y obtenga acceso a los recursos de la misma.
Además, permite reducir algunos ataques relacionados con la manipulación de
direcciones MAC y proporciona un mayor control sobre los dispositivos que
forman parte de la infraestructura de red.


\subsection{Si se desea agregar una nueva red para Alumnos en la VLAN 40 con
la ID de red 192.168.40.0/24, ¿cuál sería el procedimiento para agregar la
nueva red?}

Para agregar la nueva red se deben realizar varios cambios en la
infraestructura de red.

Primero, se debe crear la VLAN 40 en los switches que necesiten transportar
tráfico perteneciente a esta nueva red:

\begin{verbatim}
Switch(config)#vlan 40
Switch(config-vlan)#name Alumnos
Switch(config-vlan)#exit
\end{verbatim}

Posteriormente, los puertos destinados a los alumnos deben configurarse como
puertos de acceso pertenecientes a la VLAN 40:

\begin{verbatim}
Switch(config)#interface fastEthernet 0/X
Switch(config-if)#switchport mode access
Switch(config-if)#switchport access vlan 40
\end{verbatim}

Si la VLAN 40 necesita atravesar un enlace troncal entre switches o entre un
switch y un router, dicho enlace debe permitir el tráfico de la VLAN 40.

En el router que proporcione el acceso a esta VLAN se debe crear una nueva
subinterfaz. Por ejemplo, si la VLAN 40 se agregara a R2:

\begin{verbatim}
R2(config)#interface g0/1.40
R2(config-subif)#encapsulation dot1Q 40
R2(config-subif)#ip address 192.168.40.1 255.255.255.0
R2(config-subif)#no shutdown
\end{verbatim}

La dirección 192.168.40.1 funcionaría como \textit{default gateway} para
los dispositivos pertenecientes a esta nueva red.

Finalmente, los dispositivos de la VLAN 40 deberán configurarse con
direcciones pertenecientes a la red 192.168.40.0/24, por ejemplo:

\begin{verbatim}
IP:       192.168.40.10
Mask:     255.255.255.0
Gateway:  192.168.40.1
\end{verbatim}

Si se desea que esta nueva red pueda comunicarse con redes que se encuentran
detrás de otros routers, también será necesario agregar las rutas
correspondientes en los routers de la topología. En el caso de que la nueva
VLAN se encuentre conectada directamente a R2, R2 conocerá la red
192.168.40.0/24 como una red directamente conectada, mientras que R1 y R3
necesitarán rutas estáticas hacia dicha red a través de R2.


\section{Conclusión}

Durante esta práctica se implementó una topología de red compuesta por tres
routers y tres switches, utilizando VLANs para separar lógicamente diferentes
grupos de dispositivos dentro de la infraestructura.

Uno de los principales conceptos utilizados fue \textit{router-on-a-stick},
que permitió realizar el enrutamiento entre las diferentes VLANs mediante
subinterfaces configuradas sobre una misma interfaz física del router. Esta
configuración permitió mantener separadas las redes de Coordinadores,
Profesores y Becarios, mientras que los routers proporcionaron la
comunicación entre ellas cuando era necesario.

También se configuraron enlaces troncales entre los switches y los routers,
así como enlaces seriales entre los diferentes routers. En estos últimos fue
necesario identificar el extremo DCE y configurar el \textit{clock rate}
correspondiente para que los enlaces pudieran operar correctamente.

Otro aspecto importante de la práctica fue la implementación de
\textit{port-security} en los puertos de acceso de los switches. Esta
configuración proporciona una capa adicional de seguridad al limitar los
dispositivos que pueden conectarse físicamente a determinados puertos.

Finalmente, se realizaron diferentes pruebas de conectividad para verificar
el funcionamiento de la red. Durante la configuración se presentó un
problema debido a que algunos dispositivos no tenían configurado su
\textit{default gateway}. Esto provocaba que los equipos pudieran
comunicarse con dispositivos pertenecientes a su misma red, pero no pudieran
enviar correctamente paquetes hacia otras redes. Después de configurar los
gateways correspondientes, las pruebas de conectividad se realizaron
correctamente.

La práctica permitió comprender de manera más completa la relación entre
VLANs, direccionamiento IP, enlaces troncales, enrutamiento estático y
configuración de routers y switches, además de mostrar la importancia de
verificar sistemáticamente cada componente de una topología cuando existe un
problema de conectividad.

\end{document}